\documentclass[11pt]{article}
\usepackage[utf8]{inputenc}
\usepackage{amsmath,amssymb}
\usepackage{hyperref}
\title{Scalable Crispr Gene Editing Off Target for Large-Scale Settings}
\author{Anonymous}
\date{}
\begin{document}
\maketitle

\begin{abstract}
This paper studies Crispr Gene Editing Off Target. Grounded in a real, citable literature \cite{ref1,ref2,ref3,ref4,ref5,ref6,ref7,ref8},
we propose a focused method and report \textbf{illustrative} experiments.
All references below are real and verifiable (OpenAlex/arXiv); the reported
numbers are simulated to illustrate intended behaviour.
\end{abstract}

\section{Introduction}
Crispr Gene Editing Off Target is an active research area. This work targets a concrete gap identified from
the recent literature and contributes a method together with an evaluation protocol.

\section{Related Work (real literature)}
The following works, retrieved from public bibliographic APIs, frame our study:
\begin{itemize}
\item Ran et al. (2013) — "Genome engineering using the CRISPR-Cas9 system", Nature Protocols (cited 11823).
\item Komor et al. (2016) — "Programmable editing of a target base in genomic DNA without double-stranded DNA cleavage", Nature (cited 5533).
\item Ran et al. (2013) — "Double Nicking by RNA-Guided CRISPR Cas9 for Enhanced Genome Editing Specificity", Cell (cited 3384).
\item Kleinstiver et al. (2016) — "High-fidelity CRISPR–Cas9 nucleases with no detectable genome-wide off-target effects", Nature (cited 2723).
\item Kosicki et al. (2018) — "Repair of double-strand breaks induced by CRISPR–Cas9 leads to large deletions and complex rearrangements", Nature Biotechnology (cited 1919).
\item Frangoul et al. (2020) — "CRISPR-Cas9 Gene Editing for Sickle Cell Disease and β-Thalassemia", New England Journal of Medicine (cited 1893).
\end{itemize}

\section{Method}
We decompose the problem across shards with a communication-efficient aggregation rule that keeps overhead sublinear in scale. We instantiate this idea for Crispr Gene Editing Off Target and analyse its cost/quality trade-off.

\section{Experiments (Illustrative)}
\begin{center}
\begin{tabular}{lcc}
System & Primary Metric & Efficiency \\ \hline
Strong baseline & 0.812 & 1.00$\times$ \\
Ablation & 0.828 & 0.92$\times$ \\
\textbf{Ours} & \textbf{0.851} & \textbf{0.71$\times$}
\end{tabular}
\end{center}
\emph{Metrics simulated to illustrate the intended effect; not measured.}

\section{Conclusion}
We connected a real literature to a concrete method for Crispr Gene Editing Off Target. Future work will
validate the illustrative results with real experiments.

\begin{thebibliography}{9}
\bibitem{ref1} F. Ann Ran, Patrick D. Hsu, Jason Wright et al.. Genome engineering using the CRISPR-Cas9 system. \emph{Nature Protocols}, 2013. DOI:10.1038/nprot.2013.143
\bibitem{ref2} Alexis C. Komor, Y. Bill Kim, Michael S. Packer et al.. Programmable editing of a target base in genomic DNA without double-stranded DNA cleavage. \emph{Nature}, 2016. DOI:10.1038/nature17946
\bibitem{ref3} F. Ann Ran, Patrick D. Hsu, Chie-Yu Lin et al.. Double Nicking by RNA-Guided CRISPR Cas9 for Enhanced Genome Editing Specificity. \emph{Cell}, 2013. DOI:10.1016/j.cell.2013.08.021
\bibitem{ref4} Benjamin P. Kleinstiver, Vikram Pattanayak, Michelle S. Prew et al.. High-fidelity CRISPR–Cas9 nucleases with no detectable genome-wide off-target effects. \emph{Nature}, 2016. DOI:10.1038/nature16526
\bibitem{ref5} Michael Kosicki, Kärt Tomberg, Allan Bradley. Repair of double-strand breaks induced by CRISPR–Cas9 leads to large deletions and complex rearrangements. \emph{Nature Biotechnology}, 2018. DOI:10.1038/nbt.4192
\bibitem{ref6} Haydar Frangoul, David Altshuler, Maria Domenica Cappellini et al.. CRISPR-Cas9 Gene Editing for Sickle Cell Disease and β-Thalassemia. \emph{New England Journal of Medicine}, 2020. DOI:10.1056/nejmoa2031054
\bibitem{ref7} Benjamin P. Kleinstiver, Michelle S. Prew, Shengdar Q. Tsai et al.. Engineered CRISPR-Cas9 nucleases with altered PAM specificities. \emph{Nature}, 2015. DOI:10.1038/nature14592
\bibitem{ref8} Johnny H. Hu, Shannon M. Miller, Maarten H. Geurts et al.. Evolved Cas9 variants with broad PAM compatibility and high DNA specificity. \emph{Nature}, 2018. DOI:10.1038/nature26155
\end{thebibliography}
\end{document}
